% Source: https://artofproblemsolving.com/community/c3h2949414_2021_emcc_guts_round_p1_p12_exeter_math_club_competition
% Stable identity: aops-topic-2949414-q10
If endpoints are on opposite vertical sides, write them $(0,a),(6,b)$. The area below the line is $3(a+b)$, so one region has area 12 exactly when $a+b=4$ or 8. Thus $a,b$ have the same parity, and their midpoint $(3,\dfrac{a+b}{2})$ is another lattice point. The horizontal-side case is identical. If the endpoints are on adjacent sides, the corner triangle has area 12 (its area cannot exceed 18, so cannot be the complementary 24). The two integer legs have product 24 and are at most 6, hence are 4 and 6. Again both coordinate differences are even, so the midpoint is an additional lattice point. All cases contradict the condition, giving $\boxed0$.
\begin{asy}
usepackage("amsmath");
size(7cm);
draw((0,0)--(6,0)--(6,6)--(0,6)--cycle);draw((0,1)--(6,3),blue);dot((3,2),red);label("extra lattice midpoint",(3,2),N);label("area 12",(3,.5));
\end{asy}
